\chapter{Conclusion and Future Work}
\label{ch:conclusion}

\section{Conclusion}

An LLM-based repair loop produces useful information on every iteration and
then discards almost all of it. Every validation reveals which test rejected
which candidate and what that test cost; every generation reveals which repair
idea has already been tried. Current systems keep neither in a form the next
iteration can use: schedulers rank tests by a count of past rejections and
ignore how long a test takes, and prompts either carry nothing or carry a
transcript that grows without bound.

This report described \caster, which keeps that information in one small record
per test case --- how often the test has run, how often it rejected a
candidate, and how much CPU those runs took --- and spends it twice.
\textbf{Cost-aware scheduling} ranks tests by an estimated rejection rate
divided by an estimated cost, so an unpromising patch is rejected as cheaply as
possible rather than merely in as few steps as possible. \textbf{Selected
hints} replace the accumulated transcript with a small, freshly chosen set of
concrete requirements, so the prompt stays a constant size while its content
tracks what the session has just learned. The two share one record by
construction, which is what makes them a system rather than two features.

The whole mechanism is three numbers per test case, updated in constant time.
Nothing is trained, no profiling pass is needed, and neither half adds a
program execution or a model call.

\paragraph{What the pilot showed.} A pilot of $97$ ConDefects Python tasks,
three seeds each, twenty rounds per session, measured the scheduling half.

\begin{itemize}[leftmargin=*,itemsep=4pt]
\item Scheduling removed $\mathbf{60.8\,\%}$ of validation CPU against the
order the benchmark ships, interval $[-73.3, -36.3]$, while repairing exactly
the same bugs on exactly the same rounds.
\item Almost all of that came from \textbf{having a memory at all}. Any
configuration that remembered something landed within a few percent of the
others, and all of them roughly $60\,\%$ below the incumbent.
\item The new part of the rule --- \textbf{dividing by cost} --- was $2.9\,\%$
cheaper than \frtp\ over a whole session, interval $[-11.0, +4.9]$, and that
single figure hides the shape of the result. A rejection count is informative
after one observation; a cost estimate is a mean. Over rounds $1$--$5$ the
scheduler is $2.1\,\%$ \emph{more} expensive for dividing by a number it does
not yet know; over rounds $11$--$15$ it is $10.6\,\%$ cheaper; the running
total crosses over at round~$10$ and is still widening at round~$20$. The win
also concentrates in the third of tasks whose test costs actually differ. It
pays $13\,\%$ more executions for $14\,\%$ cheaper executions --- the trade it
is designed to make, and one that is invisible to a study reported in
executions alone.
\end{itemize}

The generation half has not been run. The report is therefore a measured result
for the scheduler and a declared protocol for the hints.

\section{Future Work}
\label{sec:future}

\begin{enumerate}[leftmargin=*,itemsep=8pt]

\item \textbf{Run the full corpus, and run it for longer than twenty rounds.}
The pilot screens; the corpus measures. Two of its open questions --- whether
the cost term beats \frtp, and whether a rate beats a count --- have point
estimates in the right direction and intervals that straddle zero, which is
what a corpus-scale run exists to resolve. The budget matters as much as the
corpus size here: the scheduler only overtakes \frtp\ at round~$10$ and the
gap is still growing at round~$20$, so a twenty-round budget measures the cost
term during the half of the session in which it is still learning. A longer
budget is the cheaper of the two experiments and should come first.

\item \textbf{Run Block~B.} The generation half is implemented and unit-tested
and has not been measured. It is the half the duplicate-rate literature says is
worth the most: \citet{chen2025conversational} measured $65\,\%$ of a
conversational repair system's patches to be duplicates.

\item \textbf{Add a benchmark with method-level test suites, such as
Defects4J~\citep{just2014defects4j}.} Every test in ConDefects is a
whole-program execution, so nothing can be selected away and prioritisation is
the only lever. A project-scale Java benchmark has a call graph and therefore
has selection, which is where \citet{lou2024aprregression} find selection to
dominate prioritisation. Two questions follow: does a cost-aware schedule still
help once the cheap tests have been selected away, and is the right design a
scheduler layered on a selector? That setting also tests the method's own
failure condition directly, because method-level suites are more uniform in
cost than contest pools.

\item \textbf{Add a second proposer model.} Everything the store learns is a
statistic over the candidates one model produced. If different models make
different mistakes, different tests are discriminating for each, and the
scheduler re-learns per model --- which is fine, since it learns online, but it
would mean the reported effect size is one model's. If instead the same tests
are discriminating across models, the rejection structure is a property of the
\emph{task}, which is a stronger claim: a rejection table could then be
computed once per benchmark and reused, which is also what item~$1$ needs.

\end{enumerate}
